\section{The remaining optimization problem}
\label{sec:strategy}

The preferred construction in this report uses forty-one total
coordinates, including the input. Its degree is bounded by
$2^{9566673}$, and its arithmetic interpretation remains conditional
on the named source inputs and Sun's simultaneous prime-pattern theorem.
The forty-three-coordinate semilocal alternative has the substantially
smaller degree bound $2^{236}$. For degree four, the supplied counts
are 1315 and 1335, the latter with coefficient height less than
$2^{1020}$. These are constructions and upper bounds; no optimality
claim follows from the comparisons.

\subsection{What a reduction to ten coordinates would have to accomplish}

The forty retained coordinates of the preferred source consist of the
input, eighteen components for nine elliptic abscissas, one common
slope, and twenty parameters in ten pole predicates. The finite-algebra
compiler adds one final witness. It has already removed the ordinate
roots, pole-conic roots, and the arithmetic evaluation coordinates.
The direct Lucas test removes the remaining two retained proxy
coordinates and their third, algebraic proxy. Replacing this compiler
by another primitive-element argument would leave the same retained
base and would not, by itself, decrease the count.

The geometric results above explain why each independent elliptic
point cannot simply be encoded by one rational parameter and a finite
algebraic choice. The specific cyclic point group is dense in its
two-dimensional Weil restriction. The denominator results also exclude
compressing independent parity conditions by a rank-deficient list of
monomials, or replacing the entire integrality condition by finitely
many fixed rational pullbacks of perfect-denominator predicates.
These are restrictions on particular methods, not general lower bounds
of eighteen or forty coordinates.

The most promising remaining arithmetic question is whether several
point choices or pole witnesses can be tied together while preserving
simultaneous rational solvability. A proposed sharing must prove this
simultaneous assertion. Separate existence for each pole conic, or
local solvability at every prime for each one independently, does not
provide a common rational parameter. Another possibility is an integer
relation with a different height argument requiring fewer independently
chosen elliptic labels. The direct-polynomial congruence lemma supplies
a precise interface for testing such relations, including the condition
$k>12\deg F$ for its displayed height estimate.

No construction with two through ten total coordinates has been proved
here. Two total coordinates are impossible, even for arbitrary finite
systems. The unrestricted cases of three through ten coordinates remain
unresolved by this work. For quadratic systems the proved restrictions
are stronger: at least five total coordinates and at least three
independent quadratic residuals are necessary. A successful quadratic
system in five through ten coordinates would immediately give a
quartic in exactly the same coordinates over $\mathbb Q$.

\subsection{Degree and coefficients at a fixed coordinate count}

The fixed elliptic exclusion product dominates the enormous degree of
the forty-one-coordinate construction. A better explicit height
constant, or a different size condition costing no new coordinates,
would be more consequential than a small improvement to its finite
algebra. The semilocal alternative illustrates this trade explicitly.
The remaining finite ranks are $2^{28}$ and $2^{29}$; their determinant
and separating-parameter bounds are conservative, and sparse formulas
could improve degree and coefficient bounds without saving coordinates.

For quartics, the preferred 1315-coordinate version also uses an
invertible affine change of its three private anchor witnesses, giving
exact primitive coefficient height $64\cdot3^{2^{28}}$. Removing the
last constant-chain square and inlining its first three constants gives
the 1335-coordinate alternative at height less than $2^{1020}$.
The direct field-valued Lucas test has reduced both the
number of tests and the lengths of their shared symmetric-product
computations. Further useful computation should operate on the actual
quadratic residual system and preserve an exact reconstruction
certificate. Changes of residual basis and invertible coordinate
changes may also reduce coefficient height. All comparisons must use
primitive integer coefficients after clearing numerical denominators.
Computing a huge fixed integer with extra squaring coordinates remains
a distinct trade, which is why both quartic tables are supplied.

\subsection{The arithmetic hypotheses still requiring independent review}

The most consequential review remains the source manuscript's
restricted $2$-converse, its fixed-field coefficient transfer, and
the complete five-contact height argument with its explicit constant.
The shared-point and direct-Lucas constructions additionally use
Sun's Theorem~4.1. The stated theorem and its order-compatibility
consequences were checked; the complete analytic proof of its
prime-pattern estimates has not been independently reproduced here.
The distinction between the unconditional Artin-symbol theorem of
Frei--Tinkov\'a and their conditional primitive-root theorem matters
\cite{FreiTinkova2025}. Sun's further large-prime sieve and its use
of polynomial-prime transference \cite{TaoZiegler2018} remain external
inputs, not conclusions of the polynomial checkers.

This extra input has a specific role. A common comparison index must
satisfy simultaneous congruences at different selected primes. Their
compatibility requires
\[
                    \gcd(d_{p,-},d_{q,+})=4
                         \qquad\text{for all }p,q.
\]
An odd common divisor would make the residues zero and four
incompatible. One cannot justify the common index by reusing only the
original private-witness existence statements. The conditional proofs
in this report retain the stronger simultaneous hypothesis explicitly.

\subsection{Number fields and the limits of a uniform formula}

The explicit $\mathbb Q(i)$ sections and the rejected integral contact
inputs show that the literal newer polynomials do not extend unchanged
to rings of integers. Translating the input square removes the finite
contact obstruction; a field-adapted division norm repairs conjunction
without adding variables. These are useful repairs, but neither restores
the fixed cyclic elliptic-group argument. The totally real quartic
example gives a concrete point outside that group satisfying every
finite anchor exclusion.

The general constructibility theorem further shows that no one fixed
finite existential formula can define the rings of integers in all
number-field extensions. The field-dependent existence question remains
open here. Under the claimed rational definition, an existential
definition of the embedded rational field would permit the integral
basis construction in Theorem~\ref{thm:nf-equivalence}. Undecidability
or MRDP over the ring of integers does not supply that rational-field
isolation. The next arithmetic work should address this missing
assertion or develop a different source adapted to each field.
